\documentclass[10pt,a4paper]{amsart}
\usepackage{amsmath,amssymb,amsthm,graphicx}
\usepackage{verbatim}
\newenvironment{reference}{\comment}{\endcomment}
%\newenvironment{reference}{}{}
\newenvironment{slogan}{\comment}{\endcomment}
\newenvironment{history}{\comment}{\endcomment}

% For commutative diagrams we use Xy-pic
\usepackage[all]{xy}

% We use 2cell for 2-commutative diagrams.
\xyoption{2cell}
% Theorem environments.
%
\theoremstyle{plain}
\newtheorem{theorem}[subsection]{Theorem}
\newtheorem{proposition}[subsection]{Proposition}
\newtheorem{lemma}[subsection]{Lemma}

\theoremstyle{definition}
\newtheorem{definition}[subsection]{Definition}
\newtheorem{example}[subsection]{Example}
\newtheorem{exercise}[subsection]{Exercise}
\newtheorem{situation}[subsection]{Situation}

\theoremstyle{remark}
\newtheorem{remark}[subsection]{Remark}
\newtheorem{remarks}[subsection]{Remarks}

\numberwithin{equation}{subsection}

% Macros
%
\def\lim{\mathop{\mathrm{lim}}\nolimits}
\def\colim{\mathop{\mathrm{colim}}\nolimits}
\def\Spec{\mathop{\mathrm{Spec}}}
\def\Hom{\mathop{\mathrm{Hom}}\nolimits}
\def\Ext{\mathop{\mathrm{Ext}}\nolimits}
\def\SheafHom{\mathop{\mathcal{H}\!\mathit{om}}\nolimits}
\def\SheafExt{\mathop{\mathcal{E}\!\mathit{xt}}\nolimits}
\def\Sch{\mathit{Sch}}
\def\Mor{\mathop{\mathrm{Mor}}\nolimits}
\def\Ob{\mathop{\mathrm{Ob}}\nolimits}
\def\Sh{\mathop{\mathit{Sh}}\nolimits}
\def\NL{\mathop{N\!L}\nolimits}
\def\CH{\mathop{\mathrm{CH}}\nolimits}
\def\proetale{{pro\text{-}\acute{e}tale}}
\def\etale{{\acute{e}tale}}
\def\QCoh{\mathit{QCoh}}
\def\Ker{\mathop{\mathrm{Ker}}}
\def\Im{\mathop{\mathrm{Im}}}
\def\Coker{\mathop{\mathrm{Coker}}}
\def\Coim{\mathop{\mathrm{Coim}}}

% Boxtimes
%
\DeclareMathSymbol{\boxtimes}{\mathbin}{AMSa}{"02}

%
% Macros for moduli stacks/spaces
%
\def\QCohstack{\mathcal{QC}\!\mathit{oh}}
\def\Cohstack{\mathcal{C}\!\mathit{oh}}
\def\Spacesstack{\mathcal{S}\!\mathit{paces}}
\def\Quotfunctor{\mathrm{Quot}}
\def\Hilbfunctor{\mathrm{Hilb}}
\def\Curvesstack{\mathcal{C}\!\mathit{urves}}
\def\Polarizedstack{\mathcal{P}\!\mathit{olarized}}
\def\Complexesstack{\mathcal{C}\!\mathit{omplexes}}
% \Pic is the operator that assigns to X its picard group, usage \Pic(X)
% \Picardstack_{X/B} denotes the Picard stack of X over B
% \Picardfunctor_{X/B} denotes the Picard functor of X over B
\def\Pic{\mathop{\mathrm{Pic}}\nolimits}
\def\Picardstack{\mathcal{P}\!\mathit{ic}}
\def\Picardfunctor{\mathrm{Pic}}
\def\Deformationcategory{\mathcal{D}\!\mathit{ef}}
\usepackage{cleveref}
\crefname{originallemma}{original lemma}{original lemmas}
\Crefname{originallemma}{Original lemma}{Original lemmas}
\makeatletter
\newlabel{cotangent-lemma-perfect-zero}{{0G60}{source}{}{sourcepointer-1}{}}
\newlabel{cotangent-lemma-perfect-zero@cref}{{[originallemma][0][]0G60}{[1][0][]source}{}{}{}}
\makeatother
\crefname{originallemma}{original lemma}{original lemmas}
\Crefname{originallemma}{Original lemma}{Original lemmas}
\makeatletter
\newlabel{cotangent-lemma-relation-with-naive-cotangent-complex}{{08RB}{source}{}{sourcepointer-2}{}}
\newlabel{cotangent-lemma-relation-with-naive-cotangent-complex@cref}{{[originallemma][0][]08RB}{[1][0][]source}{}{}{}}
\makeatother
\crefname{originallemma}{original lemma}{original lemmas}
\Crefname{originallemma}{Original lemma}{Original lemmas}
\makeatletter
\newlabel{cotangent-lemma-when-projective}{{08R5}{source}{}{sourcepointer-3}{}}
\newlabel{cotangent-lemma-when-projective@cref}{{[originallemma][0][]08R5}{[1][0][]source}{}{}{}}
\makeatother
\crefname{originalsection}{original section}{original sections}
\Crefname{originalsection}{Original section}{Original sections}
\makeatletter
\newlabel{cotangent-section-triangle}{{08QR}{source}{}{sourcepointer-4}{}}
\newlabel{cotangent-section-triangle@cref}{{[originalsection][0][]08QR}{[1][0][]source}{}{}{}}
\makeatother
\crefname{originallemma}{original lemma}{original lemmas}
\Crefname{originallemma}{Original lemma}{Original lemmas}
\makeatletter
\newlabel{more-morphisms-lemma-NL-formally-smooth}{{0D0L}{source}{}{sourcepointer-5}{}}
\newlabel{more-morphisms-lemma-NL-formally-smooth@cref}{{[originallemma][0][]0D0L}{[1][0][]source}{}{}{}}
\makeatother
\crefname{originallemma}{original lemma}{original lemmas}
\Crefname{originallemma}{Original lemma}{Original lemmas}
\makeatletter
\newlabel{more-morphisms-lemma-formally-unramified-differentials}{{02H9}{source}{}{sourcepointer-6}{}}
\newlabel{more-morphisms-lemma-formally-unramified-differentials@cref}{{[originallemma][0][]02H9}{[1][0][]source}{}{}{}}
\makeatother
\title{A -> A\_perf for A = F\_p[T] is flat and formally unramified but not formally etale}
\author{The Stacks Project Authors}
\date{stacks-project:examples.tex:local-snapshot-20261012; GNU FDL 1.2 or later}
\begin{document}
\maketitle
\noindent\textit{Editorial note.} All delivered mathematical text is the author's own, copied from the pinned Stacks Project chapter file examples.tex; no statement or proof step was rewritten, reordered or removed.\par\medskip
\noindent\textit{Editorial note.} The counted proof is the author's complete argument: Author source typo kept verbatim under content zero-change: '\$L\_\{A/\textbackslash{}mathbb\{F\}\_P\}\$' (capital P) occurs once in the cotangent triangle sentence; also one sentence ends with the citation and no period. No rewrite.\par\medskip
\noindent\textit{Editorial note.} The label lemma-perfect-closure-polynomial-ring names the counted statement, and the proof environment immediately following it in the source is delivered verbatim as its proof.\par\medskip
\noindent\textit{Editorial note.} The counted statement is carried by the author's own lemma environment in the chapter source.\par\medskip
\noindent\textit{Allowed standard prerequisites (editorial).}\begin{itemize}
\item Perfect closure as a colimit of Frobenius towers; faithful flatness of Frobenius on F\_p[T] via the free basis \{1, T, ..., T\textasciicircum{}(p-1)\}
\item Formal unramifiedness via vanishing of Omega (More on Morphisms, formally-unramified-differentials)
\item The cotangent complex distinguished triangle (Cotangent, section-triangle) and the naive cotangent criterion for formal smoothness (More on Morphisms, NL-formally-smooth; Cotangent, relation-with-naive-cotangent-complex)
\end{itemize}\medskip
\ifdefined\hypertarget\hypertarget{sourcepointer-1}{}\fi
\noindent\textit{Source reference.} Original lemma 0G60: sources/stacks-chow/tags.txt, source lines 20180--20180. Prerequisite pointer invoked at 3458; it names a standard fact of a sibling chapter that the counted argument uses, and every step of the counted proof built on it is retained verbatim in the delivered body.\par\medskip
\ifdefined\hypertarget\hypertarget{sourcepointer-2}{}\fi
\noindent\textit{Source reference.} Original lemma 08RB: sources/stacks-chow/tags.txt, source lines 11089--11089. Prerequisite pointer invoked at 3473; it names a standard fact of a sibling chapter that the counted argument uses, and every step of the counted proof built on it is retained verbatim in the delivered body.\par\medskip
\ifdefined\hypertarget\hypertarget{sourcepointer-3}{}\fi
\noindent\textit{Source reference.} Original lemma 08R5: sources/stacks-chow/tags.txt, source lines 11083--11083. Prerequisite pointer invoked at 3456; it names a standard fact of a sibling chapter that the counted argument uses, and every step of the counted proof built on it is retained verbatim in the delivered body.\par\medskip
\ifdefined\hypertarget\hypertarget{sourcepointer-4}{}\fi
\noindent\textit{Source reference.} Original section 08QR: sources/stacks-chow/tags.txt, source lines 11069--11069. Prerequisite pointer invoked at 3466; it names a standard fact of a sibling chapter that the counted argument uses, and every step of the counted proof built on it is retained verbatim in the delivered body.\par\medskip
\ifdefined\hypertarget\hypertarget{sourcepointer-5}{}\fi
\noindent\textit{Source reference.} Original lemma 0D0L: sources/stacks-chow/tags.txt, source lines 16316--16316. Prerequisite pointer invoked at 3471; it names a standard fact of a sibling chapter that the counted argument uses, and every step of the counted proof built on it is retained verbatim in the delivered body.\par\medskip
\ifdefined\hypertarget\hypertarget{sourcepointer-6}{}\fi
\noindent\textit{Source reference.} Original lemma 02H9: sources/stacks-chow/tags.txt, source lines 3256--3256. Prerequisite pointer invoked at 3448; it names a standard fact of a sibling chapter that the counted argument uses, and every step of the counted proof built on it is retained verbatim in the delivered body.\par\medskip
\begin{lemma}
\label{lemma-perfect-closure-polynomial-ring}
Let $A = \mathbb{F}_p[T]$ be the polynomial ring in one variable over
$\mathbb{F}_p$. Let $A_{perf}$  denote the perfect closure of $A$.
Then $A \rightarrow A_{perf}$ is flat and formally unramified,
but not formally \'etale.
\end{lemma}

\begin{proof}
Note that under the Frobenius map $F_A : A \to A$, the target copy of $A$
is a free-module over the domain with basis $\{1, T, \dots, T^{p - 1}\}$.
Thus, $F_A$ is faithfully flat, and consequently,  so is
$A \to A_{perf}$ since it is a colimit of faithfully flat maps.
Since $A_{perf}$ is a perfect ring, the relative Frobenius
$F_{A_{perf}/A}$ is a surjection. In  other words,
$A_{perf} = A[A_{perf}^p]$, which readily implies
$\Omega_{A_{perf}/A} = 0$. Then
$A \rightarrow A_{perf}$ is formally unramified by
More on Morphisms, Lemma
\ref{more-morphisms-lemma-formally-unramified-differentials}

\medskip\noindent
It suffices to show that $A \rightarrow A_{perf}$
is not formally smooth. Note that since $A$ is a
smooth $\mathbb{F}_p$-algebra, the cotangent complex
$L_{A/\mathbb{F}_P} \simeq \Omega_{A/\mathbb{F}_p}[0]$
is concentrated in degree $0$, see
Cotangent, Lemma \ref{cotangent-lemma-when-projective}. Moreover,
$L_{A_{perf}/\mathbb{F}_p} = 0$ in $D(A_{perf})$
by Cotangent, Lemma \ref{cotangent-lemma-perfect-zero}.
Consider the distinguished triangle of cotangent complexes
$$
L_{A/\mathbb{F}_p} \otimes_A A_{perf} \to
L_{A_{perf}/\mathbb{F}_p} \to
L_{A_{perf}/A} \to
(L_{A/\mathbb{F}_p} \otimes_A A_{perf})[1]
$$
in $D(A_{perf})$, see Cotangent, Section \ref{cotangent-section-triangle}.
We find $L_{A_{perf}/A} = \Omega_{A/\mathbb{F}_p} \otimes_A A_{perf}[1]$,
that is, $L_{A_{perf}/A}$ is equal to a free rank $1$ $A_{perf}$
module placed in degree $-1$. Thus $A \rightarrow A_{perf}$
is not formally smooth by
More on Morphisms, Lemma \ref{more-morphisms-lemma-NL-formally-smooth}
and
Cotangent, Lemma \ref{cotangent-lemma-relation-with-naive-cotangent-complex}.
\end{proof}

\end{document}
